\documentclass[a4paper,11pt,fleqn]{scrartcl}
\usepackage{../../Styles/Eval}
\usepackage{longtable}

\newcommand{\chnum}{Chapitre 21 et 14}
\newcommand{\chtitle}{Corrigé -- Dipôle RC et Mouvement d'un fluide}
\newcommand{\dtype}{DS 5}

\begin{document}
\maketitle

\section{Installation sanitaire}

\begin{enumerate}
    \item L'eau est considérée comme un fluide incompressible en régime permanent.\\
    Le débit volumique se conserve donc tout au long de l'installation :
    \[
    D_1 = D_2
    \]

    \item Or :
    \[
    D = v \times S
    \]

    Donc :
    \[
    v_1S_1 = v_2S_2
    \]

    D'où :
    \[
    v_2 = \dfrac{v_1S_1}{S_2}
    \]

    Application numérique :
    \[
    v_2 = \dfrac{1,0 \times 2,0\times10^{-4}}{1,1\times10^{-4}}
    \]

    \[
    v_2 \approx \SI{1,8}{\meter\per\second}
    \]

    \item La relation de Bernoulli entre les points 1 et 2 s'écrit :
    \[
    P_1 + \dfrac{1}{2}\rho v_1^2 + \rho gz_1
    =
    P_2 + \dfrac{1}{2}\rho v_2^2 + \rho gz_2
    \]

    On isole $P_2$ :
    \[
    P_2 =
    P_1 + \dfrac{1}{2}\rho(v_1^2-v_2^2)+\rho g(z_1-z_2)
    \]

    Application numérique :
    \[
    P_2 =
    3,0\times10^5
    +\dfrac{1}{2}\times1000\times(1,0^2-1,8^2)
    +1000\times9,81\times(0-5)
    \]

    \[
    P_2 \approx 2,4\times10^5\ \si{\pascal}
    \]

    Conversion en bar :
    \[
    P_2 \approx \SI{2,4}{\bar}
    \]

    \item Un surpresseur est nécessaire uniquement si la pression est inférieure à \SI{2,0}{\bar}.\\

    Or :
    \[
    P_2 \approx \SI{2,4}{\bar}
    \]

    Donc la pression reste suffisante :\\
    \textbf{un surpresseur n'est pas nécessaire.}
\end{enumerate}

\section{Capteur capacitif}

\subsection{Étude théorique de la charge d'un dipôle RC}

\begin{enumerate}
    \item Loi des mailles :
    \[
    E-u_R-u_C=0
    \]

    soit :
    \[
    E=u_R+u_C
    \]

    \item Loi d'Ohm :
    \[
    u_R=Ri
    \]

    \item Relation pour le condensateur :
    \[
    i=C\dfrac{du_C}{dt}
    \]

    \item On remplace $u_R$ dans la loi des mailles :
    \[
    E=Ri+u_C
    \]

    Puis :
    \[
    E=RC\dfrac{du_C}{dt}+u_C
    \]

    Donc :
    \[
    \dfrac{du_C}{dt}+\dfrac{u_C}{RC}
    =
    \dfrac{E}{RC}
    \]

    \item On considère :
    \[
    u_C=E\left(1-e^{-t/RC}\right)
    \]

    Calculons sa dérivée :
    \[
    \dfrac{du_C}{dt}
    =
    \dfrac{E}{RC}e^{-t/RC}
    \]

    Alors :
    \[
    \dfrac{du_C}{dt}+\dfrac{u_C}{RC}
    =
    \dfrac{E}{RC}e^{-t/RC}
    +
    \dfrac{E(1-e^{-t/RC})}{RC}
    \]

    \[
    =
    \dfrac{E}{RC}
    \]

    La fonction proposée vérifie bien l'équation différentielle.
\end{enumerate}

\subsection{Étude expérimentale de la charge d'un dipôle RC}

\begin{enumerate}[resume]
    \item Graphiquement, le temps caractéristique correspond à :
    \[
    u_C = 0,63\times U_f
    \]

    La lecture graphique donne :
    \[
    \tau \approx \SI{3,3}{ms}
    \]

    \item La relation entre le temps caractéristique, la résistance et la capacité est :
    \[
    \tau = RC
    \]

    \item On en déduit :
    \[
    C=\dfrac{\tau}{R}
    \]

    Application numérique :
    \[
    C=\dfrac{3,3\times10^{-3}}{330}
    \]

    \[
    C=1,0\times10^{-5}\ \si{\farad}
    \]

    soit :
    \[
    C=\SI{10}{\micro\farad}
    \]
\end{enumerate}

\subsection{Étude d'un condensateur à capacité variable}

\begin{enumerate}[resume]
    \item On part de :
    \[
    C=\dfrac{\varepsilon S}{d}
    \]

    Donc :
    \[
    \varepsilon = \dfrac{Cd}{S}
    \]

    Unités :
    \[
    [\varepsilon]=\dfrac{\si{\farad}\times\si{\meter}}{\si{\meter\squared}}
    \]

    \[
    [\varepsilon]=\si{\farad\per\meter}
    \]

    \item Lorsque l'on rapproche les armatures, la distance $d$ diminue.\\

    Or :
    \[
    C=\dfrac{\varepsilon S}{d}
    \]

    Donc si $d$ diminue, la capacité $C$ augmente.
\end{enumerate}

\subsection{Réalisation d'un capteur de position}

\begin{enumerate}[resume]
    \item Les deux plaques métalliques constituent les armatures d'un condensateur et la colle joue le rôle d'isolant.\\

    La capacité dépend de l'épaisseur de colle :
    \[
    C=\dfrac{\varepsilon S}{d}
    \]

    Si l'épaisseur de colle augmente, la distance $d$ augmente donc la capacité diminue.\\
    Si l'épaisseur diminue, la capacité augmente.\\

    En mesurant la capacité du condensateur grâce au circuit RC (par exemple via le temps caractéristique $\tau=RC$), on peut déterminer l'épaisseur de colle et vérifier qu'elle est correcte.
\end{enumerate}

\end{document}